\section{Kesimpulan}

Studi ini menghasilkan model komputasional berbasis simulasi untuk
mengeksplorasi Peto's Paradox pada tingkat populasi sel. Model menggabungkan
mutasi multi-hit, repair, apoptosis, pertumbuhan kanker, kontrol imun, dan
definisi cancer escape sebagai kondisi ketika sel kanker melewati ambang imun.
Dalam 100 percobaan per skenario, skenario large long-lived naive menghasilkan
escape probability \NaiveEscapeProbability{}, sedangkan mekanisme protektif
seperti pembelahan lebih lambat, repair lebih baik, ambang mutasi lebih tinggi,
dan kontrol imun lebih kuat menurunkan peluang tersebut.

Analisis dataset Vincze dkk. digunakan sebagai konteks empiris. CMR lintas
mamalia tidak meningkat secara proporsional seiring massa tubuh atau exposure
proxy. Slope log massa tubuh terhadap log CMR bernilai \MassSlope\ dengan
$R^2=\MassRsq$, sedangkan exposure proxy memberi $R^2=\ExposureRsq$. Hasil ini
mendukung motivasi simulasi: ukuran tubuh dan umur saja belum memadai untuk
menjelaskan variasi risiko kanker antarspesies.

Kontribusi utama proyek ini adalah kerangka simulasi eksploratif yang sederhana,
transparan, dan reproduksibel untuk menguji intuisi mekanistik Peto's Paradox. Model tidak
membuktikan mekanisme spesies tertentu dan tidak memprediksi risiko absolut.
Namun, hasilnya menunjukkan bahwa perubahan kecil pada proses seluler dan imun
dapat menurunkan cancer escape pada kondisi ukuran dan umur yang lebih besar.
Pengembangan berikutnya perlu memasukkan parameter jaringan yang lebih realistis,
kontrol filogenetik, dan data molekuler agar simulasi dapat bergerak dari
eksplorasi konseptual menuju pengujian biologis yang lebih spesifik.
